% 第8章：极分解（伸缩与正交转向的分离、可逆方阵、奇异方阵与矩形矩阵的极分解、矩阵函数的谱表示）。

\chapter{极分解}
\label{sec:ls-polar-decomposition}
\label{chap:polar-decomposition}

非零复数$z$可以写成$z=e^{\mathrm i\theta}r$：先乘正数$r$改变长度，再乘单位复数$e^{\mathrm i\theta}$转动方向。本章对矩阵提出同样的问题：给定$\bm A$，能否把它写成$\bm A=\bm W\bm H$，使$\bm H$沿一组互相垂直的方向作非负伸缩，$\bm W$保持长度？这样的分解是否唯一？先转向、后伸缩是否也可以？矩阵奇异或不是方阵时，又能保留多少？

第~\ref{chap:singular-value-decomposition}章的奇异值分解已经把$\bm A$写成“转向、伸缩、再转向”的三段。本章把三段重新组合成两段，并关心另一个问题：两个因子能否直接由$\bm A$求出，而不先求出全部奇异向量。

第~\ref{sec:ls-polar-square-root-equation}节先讨论伸缩与正交转向怎样分离：利用正交因子保持长度把它消去，得到只含伸缩因子的开方方程，沿谱开方构造唯一的半正定伸缩因子，并据此定义极分解。伸缩因子在任何情形下都这样求出，正交因子的求法则随矩阵的情形而变，后面三节各处理一种情形：第~\ref{sec:ls-polar-orthogonal-factor}节处理可逆方阵，撤销伸缩即得唯一的正交因子，并比较先伸缩后转向与先转向后伸缩两种写法；第~\ref{sec:ls-polar-singular-best-orthogonal}节处理奇异方阵，借助奇异值分解构造正交因子，并说明核怎样使它失去唯一性；第~\ref{sec:ls-polar-rectangular}节处理矩形矩阵，说明保长因子只能是列标准正交或行标准正交的矩阵。

第~\ref{sec:ls-matrix-functions}节把开方和求逆中共同使用的“沿谱改写特征值”归纳为对称矩阵的函数。

本章用到第~\ref{chap:eigen-spectral-decomposition}章的谱定理、谱投影与定号性判据，奇异情形还用到第~\ref{chap:singular-value-decomposition}章的奇异值分解。极分解的正交因子还有一个最优性：第~\ref{chap:matrix-approximation}章将证明，在Frobenius范数下它是离$\bm A$最近的正交矩阵（定理~\ref{thm:ls-best-orthogonal-approximation}）。这需要比较全部正交矩阵，属于约束下的矩阵逼近问题，本章不展开。

\section{伸缩与正交转向的分离}
\label{sec:ls-polar-square-root-equation}

本节说明伸缩与转向怎样分开：先消去正交因子，得到只含伸缩因子的开方方程；再沿谱开方，构造唯一的半正定伸缩因子；最后据此定义极分解，并说明正交因子的求法为什么要按情形讨论。

\subsection{开方方程}
\label{sec:ls-polar-square-root-subsection}

先看一个具体矩阵。本章用以下矩阵作为贯穿例：
\[
 \bm A=\begin{bmatrix}0&-1\\3&0\end{bmatrix}.
\]
它把$\bm e_1$送到$3\bm e_2$，把$\bm e_2$送到$-\bm e_1$。这个作用既改变长度，又改变方向：$\bm e_1$被拉长为原来的$3$倍，同时两个坐标轴都被转了$90$度。我们希望把它拆成两步：先由对称正定的$\bm H$沿一组正交方向伸缩，再由正交矩阵$\bm W$保长地改变方向，即求$\bm A=\bm W\bm H$。

本节的推导对一般方阵$\bm A\in\mathbb R^{d\times d}$成立，例子取可逆矩阵。$\bm A$可逆时，它不把任何非零向量压成零，伸缩可以撤销，所求的$\bm H$应当正定；$\bm A$奇异时，有些方向被压成零，$\bm H$只能半正定。正交矩阵既可以是旋转，也可以包含反射，这里不加限制。矩形矩阵的情形在第~\ref{sec:ls-polar-rectangular}节处理。

两个因子同时未知，先设法消去一个。若$\bm A=\bm W\bm H$，其中$\bm W^\top\bm W=\bm I$、$\bm H^\top=\bm H$，则把等式两边分别转置后与原式相乘，得到
\begin{equation}
 \bm A^\top\bm A
 =\bm H^\top\bm W^\top\bm W\bm H
 =\bm H^2 .
 \label{eq:ls-polar-square-root-equation}
\end{equation}
正交因子在相乘时消失了，因为它保持长度：
\[
 \lVert\bm A\bm x\rVert_2=\lVert\bm W\bm H\bm x\rVert_2=\lVert\bm H\bm x\rVert_2 .
\]

记$\bm S=\bm A^\top\bm A$，它完全由已知的$\bm A$确定。于是原问题的前一半化为一个开方问题：找一个对称正定矩阵$\bm H$，使$\bm H^2=\bm S$。需要注意，式~\eqref{eq:ls-polar-square-root-equation}只是分解若存在就必须满足的条件；这样的$\bm H$是否存在、由它能否再得到正交的$\bm W$，都还没有证明。

开方之前先确认$\bm S$的结构。它对称，因为$(\bm A^\top\bm A)^\top=\bm A^\top\bm A$；对任意$\bm x$，
\[
 \bm x^\top\bm S\bm x=(\bm A\bm x)^\top(\bm A\bm x)=\lVert\bm A\bm x\rVert_2^2\geq0,
\]
所以$\bm S$对称半正定，可以使用第~\ref{sec:ls-spectral-theory}节的正交谱分解，全部特征值非负。$\bm A$可逆时，非零$\bm x$使$\bm A\bm x\neq\bm0$，上式严格大于零，$\bm S$正定，全部特征值为正。

在贯穿例中，先算出$\bm S$：
\[
 \begin{aligned}
 \bm S=\bm A^\top\bm A
 &=\begin{bmatrix}0&3\\-1&0\end{bmatrix}
 \begin{bmatrix}0&-1\\3&0\end{bmatrix}\\
 &=\begin{bmatrix}9&0\\0&1\end{bmatrix}.
 \end{aligned}
\]
这里$\bm S$恰好是对角矩阵，最直接的做法是沿两个坐标轴分别取正平方根，得到候选$\bm H=\operatorname{diag}(3,1)$。一般的$\bm S$不是对角矩阵，下一小节用谱分解把它化到这种情形。

\subsection{伸缩因子的构造}
\label{sec:ls-matrix-square-root}

对上一小节的$\bm S=\bm A^\top\bm A$作正交谱分解
\[
 \bm S=\bm Q\bm\Lambda\bm Q^\top,
\]
其中$\bm Q=[\bm q_1,\ldots,\bm q_d]$正交，$\bm\Lambda=\operatorname{diag}(\lambda_1,\ldots,\lambda_d)$，全部$\lambda_i\geq0$。先看$\bm A$可逆、全部$\lambda_i>0$的情形。

在特征方向$\bm q_i$上，$\bm S$只作标量伸缩$\lambda_i$倍。既然希望$\bm H$作用两次等于$\bm S$，自然尝试让$\bm H$沿同一方向伸缩$\sqrt{\lambda_i}$倍，即构造
\begin{equation}
 \bm H=\bm Q\operatorname{diag}\bigl(\sqrt{\lambda_1},\ldots,\sqrt{\lambda_d}\bigr)\bm Q^\top .
 \label{eq:ls-spectral-square-root}
\end{equation}
这是我们主动选取的一个候选，并没有预先假定未知的$\bm H$必须与$\bm S$共用特征基；它是否满足要求，要靠验证。

验证只需矩阵乘法。记$\bm\Lambda^{1/2}=\operatorname{diag}(\sqrt{\lambda_1},\ldots,\sqrt{\lambda_d})$，由$\bm Q^\top\bm Q=\bm I$，
\[
 \bm H^2
 =\bm Q\bm\Lambda^{1/2}(\bm Q^\top\bm Q)\bm\Lambda^{1/2}\bm Q^\top
 =\bm Q\bm\Lambda\bm Q^\top
 =\bm S .
\]
又
\[
 \bm H^\top=\bm Q(\bm\Lambda^{1/2})^\top\bm Q^\top=\bm H ,
\]
而式~\eqref{eq:ls-spectral-square-root}本身就是$\bm H$的正交谱分解，特征值$\sqrt{\lambda_i}$全为正，所以$\bm H$对称正定。

所需的伸缩因子已经求出，记作$\bm S^{1/2}$，称为$\bm S$的正定平方根，也称\term{矩阵平方根}（matrix square root）。

构造本身只用到$\lambda_i\geq0$，所以对半正定的$\bm S$同样可行，此时得到半正定的$\bm H$。还剩一个问题：$\bm Q$不唯一，特征值有重根时，重根空间内可以换用别的标准正交基。

方程$\bm X^2=\bm S$本身也可能有许多解。例如$\bm I$、$-\bm I$与$\operatorname{diag}(1,-1)$的平方都是$\bm I_2$，后两个对称却不半正定；不对称的$\left[\begin{smallmatrix}0&1\\0&0\end{smallmatrix}\right]$与零矩阵的平方都是零矩阵。

所以“平方等于$\bm S$”这一代数条件不足以确定伸缩因子。下面的定理说明，只要要求解对称半正定，答案就唯一，因而与$\bm Q$的选取无关。

\begin{theorem}[半正定平方根的存在与唯一]
\label{thm:ls-psd-square-root}
设$\bm S\in\mathbb R^{d\times d}$对称半正定，有正交谱分解$\bm S=\bm Q\bm\Lambda\bm Q^\top$。则存在唯一的对称半正定矩阵$\bm X$满足$\bm X^2=\bm S$，它就是式~\eqref{eq:ls-spectral-square-root}给出的$\bm S^{1/2}$。$\bm S$正定时，$\bm S^{1/2}$也正定。
\end{theorem}
\begin{proof}
存在性与正定性已由上面的验证给出：$\bm S$半正定时全部$\lambda_i\geq0$，$\sqrt{\lambda_i}$有定义，式~\eqref{eq:ls-spectral-square-root}是对称半正定矩阵且平方为$\bm S$；$\bm S$正定时全部$\sqrt{\lambda_i}>0$。

下证唯一性。设$\bm X$对称半正定且$\bm X^2=\bm S$。对$\bm X$本身作正交谱分解
\[
 \bm X=\sum_{j=1}^d\mu_j\bm p_j\bm p_j^\top ,
\]
其中$\bm p_1,\ldots,\bm p_d$是标准正交基，$\mu_j\geq0$。由$\bm p_j^\top\bm p_k=\delta_{j,k}$，
\[
 \bm S=\bm X^2=\sum_{j=1}^d\mu_j^2\bm p_j\bm p_j^\top ,
\]
所以$\bm S\bm p_j=\mu_j^2\bm p_j$，每个$\bm p_j$都是$\bm S$的特征向量。

取$\bm S$的任一特征值$\lambda$，考察特征空间$\ker(\bm S-\lambda\bm I)$。把任意$\bm x$展开为$\bm x=\sum_jc_j\bm p_j$，则
\[
 (\bm S-\lambda\bm I)\bm x=\sum_{j=1}^dc_j(\mu_j^2-\lambda)\bm p_j ,
\]
由$\bm p_j$线性无关，它为零当且仅当凡$\mu_j^2\neq\lambda$的$j$都有$c_j=0$。因此$\ker(\bm S-\lambda\bm I)$恰由满足$\mu_j^2=\lambda$的那些$\bm p_j$张成。

对这些$j$，$\mu_j\geq0$迫使$\mu_j=\sqrt\lambda$，于是对特征空间中的每个$\bm x$，
\[
 \bm X\bm x=\sum_{\mu_j^2=\lambda}c_j\mu_j\bm p_j=\sqrt\lambda\,\bm x .
\]
另一方面，由式~\eqref{eq:ls-spectral-square-root}，
\[
 \bm S^{1/2}\bm q_i=\sqrt{\lambda_i}\bm q_i ,
\]
而$\ker(\bm S-\lambda\bm I)$由满足$\lambda_i=\lambda$的$\bm q_i$张成，所以$\bm S^{1/2}$在这个特征空间上同样是乘以$\sqrt\lambda$。

这样，$\bm X$与$\bm S^{1/2}$在$\bm S$的每个特征空间上作用相同。由谱定理，这些特征空间合起来张成$\mathbb R^d$（各$\bm q_i$都落在其中某一个里），线性映射在一组基上相同即处处相同，故$\bm X=\bm S^{1/2}$。
\end{proof}

唯一性的论证直接说明了两件事。第一，在重根空间内换基，相应的$\bm q_i$变了，但特征空间和其上的作用“乘以$\sqrt\lambda$”不变，所以$\bm S^{1/2}$不依赖$\bm Q$的选取。第二，唯一性只针对半正定解；去掉这一要求，$\bm X^2=\bm S$一般有许多解，$-\bm S^{1/2}$就是其中之一。

在贯穿例中，$\bm S=\operatorname{diag}(9,1)$，可取$\bm Q=\bm I$，得到$\bm H=\operatorname{diag}(3,1)$。直接核对：
\[
 \bm H^2=\operatorname{diag}(9,1)=\bm S .
\]
上一小节凭直觉写出的候选正是唯一的正定平方根。

伸缩因子已经求出，正交因子的求法随矩阵的情形而变，放到后面三节讨论。

这里所做的“把每个特征值换成它的平方根”，将在第~\ref{sec:ls-matrix-functions}节成为矩阵函数的一个具体例子；但构造和验证本身只用到谱分解、标量开方与矩阵乘法。

所得的平方根带来一个后文直接使用的长度恒等式。由$\bm S^{1/2}$对称及$(\bm S^{1/2})^2=\bm S$，对半正定$\bm S$与任意$\bm x$，
\begin{equation}
 \lVert\bm S^{1/2}\bm x\rVert_2^2
 =\bm x^\top\bm S^{1/2}\bm S^{1/2}\bm x
 =\bm x^\top\bm S\bm x .
 \label{eq:ls-square-root-length}
\end{equation}
它把第~\ref{sec:ls-weighted-inner-products}节以$\bm S$加权的长度化为普通欧氏长度。

对$\bm S=\bm A^\top\bm A$，它给出
\[
 \lVert\bm H\bm x\rVert_2=\lVert\bm A\bm x\rVert_2 ,
\]
说明$\bm H$承担了$\bm A$的全部长度变化；但这并不表示$\bm H\bm x$与$\bm A\bm x$方向相同，方向的差别正要由$\bm W$补上。

第~\ref{chap:stationary-iterations}章分析Gauss--Seidel迭代的收缩时，就用式~\eqref{eq:ls-square-root-length}把能量范数写成欧氏长度。

\subsection{极分解的定义}
\label{sec:ls-polar-definition}

上两小节说明，若$\bm A$能写成正交因子与半正定伸缩因子之积，伸缩因子只能是$(\bm A^\top\bm A)^{1/2}$。把这样的分解单独命名。
\begin{definition}[极分解]
\label{def:matrix-analysis-polar-decomposition}
设$\bm A\in\mathbb R^{d\times d}$。若$\bm W$是正交矩阵，$\bm H$是对称半正定矩阵，且
\begin{equation}
\begin{gathered}
 \bm A=\bm W\bm H\\
 \bm W^\top\bm W=\bm I,\quad \bm H=(\bm A^\top\bm A)^{1/2},
\end{gathered}
\label{eq:matrix-analysis-polar-decomposition}
\end{equation}
则称$\bm A=\bm W\bm H$为$\bm A$的\term{极分解}（polar decomposition），也称右极分解，$\bm W$称为正交因子，$\bm H$称为伸缩因子。
\end{definition}

式中$\bm H=(\bm A^\top\bm A)^{1/2}$并不是额外的要求：由式~\eqref{eq:ls-polar-square-root-equation}，任何满足$\bm A=\bm W\bm H$、$\bm W$正交、$\bm H$对称半正定的分解都有$\bm H^2=\bm A^\top\bm A$，再由定理~\ref{thm:ls-psd-square-root}，$\bm H$只能是$(\bm A^\top\bm A)^{1/2}$。所以极分解中的伸缩因子总是唯一的，问题只在正交因子。

同一个作用也可以先转向、后伸缩，即求$\bm A=\bm K\bm W'$，其中$\bm K$对称半正定、$\bm W'$正交。仿照式~\eqref{eq:ls-polar-square-root-equation}，这次从右侧消去正交因子：
\[
 \bm A\bm A^\top=\bm K\bm W'\bm W'^\top\bm K^\top=\bm K^2 .
\]
所以$\bm K$只能是$\bm A\bm A^\top$的半正定平方根，$\bm A$可逆时它正定。这样的分解称为左极分解。$\bm A\bm A^\top$与$\bm A^\top\bm A$一般不同，左右两个伸缩因子也就不同。

剩下的是正交因子怎样求。它取决于$\bm A$是否可逆、是否为方阵：可逆方阵的伸缩因子可逆，撤销伸缩即得正交因子（第~\ref{sec:ls-polar-orthogonal-factor}节）；奇异方阵有被压成零的方向，正交因子要借助奇异值分解构造，并且不唯一（第~\ref{sec:ls-polar-singular-best-orthogonal}节）；矩形矩阵的输入与输出不在同一空间，保长因子只能是列标准正交或行标准正交的矩阵（第~\ref{sec:ls-polar-rectangular}节）。

\section{可逆矩阵的极分解}
\label{sec:ls-polar-orthogonal-factor}

可逆方阵的伸缩因子正定，因而可逆，正交因子可以直接由撤销伸缩求出。整个构造分四步：先算$\bm S=\bm A^\top\bm A$；再作正交谱分解$\bm S=\bm Q\bm\Lambda\bm Q^\top$，取$\bm H=\bm Q\bm\Lambda^{1/2}\bm Q^\top$；然后取$\bm H^{-1}=\bm Q\bm\Lambda^{-1/2}\bm Q^\top$，令$\bm W=\bm A\bm H^{-1}$；最后核验$\bm W^\top\bm W=\bm I$与$\bm W\bm H=\bm A$。本节依次说明这几步、证明两个因子唯一，再给出先转向后伸缩的左极分解。

\subsection{构造过程}
\label{sec:ls-polar-invertible-construction}

回到$\bm A=\bm W\bm H$。刚求出的$\bm H$正定，因而可逆，在等式两边右乘$\bm H^{-1}$，得到正交因子唯一可能的候选
\[
 \bm W=\bm A\bm H^{-1}.
\]
右乘而不是左乘，是由原方程中因子的顺序决定的：$\bm H$先作用，要撤销它的伸缩，就要在$\bm A$的右侧作用$\bm H^{-1}$。

$\bm H^{-1}$同样可以沿特征方向写出。$\bm H$沿$\bm q_i$伸缩$\sqrt{\lambda_i}$倍，撤销这一伸缩就是沿同一方向乘以$1/\sqrt{\lambda_i}$。令
\[
 \bm H^{-1}=\bm Q\operatorname{diag}\bigl(1/\sqrt{\lambda_1},\ldots,1/\sqrt{\lambda_d}\bigr)\bm Q^\top ,
\]
由$\bm Q^\top\bm Q=\bm I$，它与$\bm H$相乘得
\[
 \bm H\bm H^{-1}=\bm Q\bm\Lambda^{1/2}\bm\Lambda^{-1/2}\bm Q^\top=\bm Q\bm I\bm Q^\top=\bm I ,
\]
确实是逆矩阵。它也对称正定，记作$\bm S^{-1/2}$，称为$\bm S$的逆平方根。

定义了$\bm W=\bm A\bm H^{-1}$之后，必须检查它满足原来要求的性质。由$(\bm H^{-1})^\top=(\bm H^\top)^{-1}=\bm H^{-1}$及$\bm A^\top\bm A=\bm S=\bm H^2$，
\[
 \bm W^\top\bm W
 =\bm H^{-1}\bm A^\top\bm A\bm H^{-1}
 =\bm H^{-1}\bm H^2\bm H^{-1}
 =\bm I ,
\]
所以$\bm W$正交。又
\[
 \bm W\bm H=\bm A\bm H^{-1}\bm H=\bm A .
\]
正交性和分解等式都成立，存在性才算证完。

现在把贯穿例算完。$\bm H=\operatorname{diag}(3,1)$，所以$\bm H^{-1}=\operatorname{diag}(1/3,1)$，右乘得
\[
 \begin{aligned}
 \bm W=\bm A\bm H^{-1}
 &=\begin{bmatrix}0&-1\\3&0\end{bmatrix}
 \begin{bmatrix}1/3&0\\0&1\end{bmatrix}\\
 &=\begin{bmatrix}0&-1\\1&0\end{bmatrix}.
 \end{aligned}
\]
逐项核对：
\[
 \bm W^\top\bm W=\bm I,\qquad
 \bm W\bm H=\begin{bmatrix}0&-1\\3&0\end{bmatrix}=\bm A .
\]
$\bm W$把$\bm e_1$送到$\bm e_2$、把$\bm e_2$送到$-\bm e_1$，是逆时针旋转$90$度。

于是原变换可以读成：先沿$\bm e_1$伸长$3$倍、沿$\bm e_2$保持不变，再整体旋转$90$度。在这个解题过程中，$\bm S$是由$\bm A$直接算出的中间量，$\bm H$是对$\bm S$开方得到的因子，$\bm W$则是撤销伸缩后剩下的因子。

\subsection{唯一性}
\label{sec:ls-polar-invertible-uniqueness}

上面的构造给出了一组因子。下面的定理说明，可逆时极分解只有这一组。
\begin{theorem}[可逆矩阵的极分解]
\label{thm:ls-polar-invertible}
设$\bm A\in\mathbb R^{d\times d}$可逆，则$\bm A$有极分解，且两个因子都唯一：
\begin{equation}
 \bm H=(\bm A^\top\bm A)^{1/2},\qquad
 \bm W=\bm A(\bm A^\top\bm A)^{-1/2},
 \label{eq:ls-polar-invertible-factors}
\end{equation}
其中$\bm H$对称正定。
\end{theorem}
\begin{proof}
存在性：由第~\ref{sec:ls-polar-square-root-equation}节，$\bm S=\bm A^\top\bm A$对称正定；由定理~\ref{thm:ls-psd-square-root}，$\bm H=\bm S^{1/2}$对称正定，从而可逆；上面已验证$\bm W=\bm A\bm H^{-1}$正交且$\bm W\bm H=\bm A$。

唯一性：设$\bm A=\bm W'\bm H'$是任一极分解。由式~\eqref{eq:ls-polar-square-root-equation}，$\bm H'^2=\bm S$，又$\bm H'$对称半正定，由定理~\ref{thm:ls-psd-square-root}得$\bm H'=\bm S^{1/2}=\bm H$。

$\bm H$可逆，在$\bm A=\bm W'\bm H$两边右乘$\bm H^{-1}$得$\bm W'=\bm A\bm H^{-1}=\bm W$。
\end{proof}

由$\det\bm A=\det\bm W\det\bm H$及$\det\bm H>0$，$\det\bm W$与$\det\bm A$同号：$\det\bm A>0$时$\bm W$是旋转，$\det\bm A<0$时$\bm W$包含一次反射。这正是第~\ref{sec:ls-polar-square-root-equation}节不限制正交作用类型的原因。

第~\ref{sec:ls-matrix-square-root}节的长度恒等式还可以与Gram因子比较。若只需表示同一个长度，平方根并不是唯一的选择：由定理~\ref{thm:ls-gram-factor}，正定$\bm S$还可写成$\bm S=\bm L^\top\bm L$，$\bm L$可逆，此时同样有
\[
 \lVert\bm L\bm x\rVert_2^2=\bm x^\top\bm S\bm x=\lVert\bm S^{1/2}\bm x\rVert_2^2 .
\]
这只说明两个向量$\bm L\bm x$与$\bm S^{1/2}\bm x$长度相同，并不表示两个矩阵对向量的作用相同：$\bm L$一般不对称，$\bm S^{1/2}$则由定理~\ref{thm:ls-psd-square-root}唯一确定。

极分解说明二者恰好相差一个正交因子。由$\bm L^\top\bm L=\bm S$，$\bm L$的极分解是
\[
 \bm L=(\bm L\bm S^{-1/2})\bm S^{1/2},
\]$\bm L$与$\bm S^{1/2}$相差正交因子$\bm L\bm S^{-1/2}$，这解释了二者为何给出相同的长度。

第~\ref{chap:direct-linear-solvers}章的Cholesky分解把$\bm L$取成三角矩阵，那时作用$\bm L^{-1}$只需解三角方程，不必求出完整的谱分解。

\subsection{左极分解}
\label{sec:ls-polar-left}

第~\ref{sec:ls-polar-definition}节说明，先转向后伸缩的伸缩因子只能是$(\bm A\bm A^\top)^{1/2}$。可逆时，它也可以由右极分解直接得到。
左极分解与右极分解用的是同一个正交因子。由右极分解$\bm A=\bm W\bm H$，在$\bm H$左侧插入$\bm W^\top\bm W=\bm I$，得
\[
 \bm A=\bm W\bm H\bm W^\top\bm W=(\bm W\bm H\bm W^\top)\bm W .
\]
令$\bm K=\bm W\bm H\bm W^\top$。它对称；由$\bm H$的谱分解，
\[
 \bm K=(\bm W\bm Q)\bm\Lambda^{1/2}(\bm W\bm Q)^\top ,
\]
而$\bm W\bm Q$正交，所以$\bm K$正定，特征值与$\bm H$相同，特征方向是$\bm H$的特征方向经$\bm W$转过后的方向。

再看$\bm K$的平方：
\[
 \bm K^2=\bm W\bm H^2\bm W^\top=\bm W\bm A^\top\bm A\bm W^\top=\bm A\bm A^\top ,
\]
最后一步用了$\bm A\bm A^\top=\bm W\bm H\bm H\bm W^\top$。由定理~\ref{thm:ls-psd-square-root}，$\bm K=(\bm A\bm A^\top)^{1/2}$。

\begin{theorem}[左极分解]
\label{thm:polar-left-invertible}
设$\bm A\in\mathbb R^{d\times d}$可逆，$\bm A=\bm W\bm H$是它的极分解。则
\begin{equation}
 \bm A=\bm K\bm W,\qquad\bm K=(\bm A\bm A^\top)^{1/2}=\bm W\bm H\bm W^\top ,
 \label{eq:ls-polar-left}
\end{equation}
其中$\bm K$对称正定。满足$\bm A=\bm K'\bm W'$、$\bm K'$对称半正定、$\bm W'$正交的分解只有这一个，称为$\bm A$的\term{左极分解}（left polar decomposition）。
\end{theorem}
\begin{proof}
存在性已在定理前验证。唯一性：若$\bm A=\bm K'\bm W'$，则$\bm K'^2=\bm A\bm A^\top$，由定理~\ref{thm:ls-psd-square-root}得$\bm K'=(\bm A\bm A^\top)^{1/2}=\bm K$。$\bm K$可逆，故
\[
 \bm W'=\bm K^{-1}\bm A=\bm K^{-1}\bm K\bm W=\bm W .
\]
\end{proof}

在贯穿例中，$\bm A\bm A^\top=\operatorname{diag}(1,9)$，所以$\bm K=\operatorname{diag}(1,3)$；也可以由$\bm K=\bm W\bm H\bm W^\top$核对：旋转$90$度把$\bm H$沿$\bm e_1$的伸长$3$倍转到沿$\bm e_2$。直接相乘，
\[
 \begin{aligned}
 \bm K\bm W&=\begin{bmatrix}1&0\\0&3\end{bmatrix}\begin{bmatrix}0&-1\\1&0\end{bmatrix}\\
 &=\begin{bmatrix}0&-1\\3&0\end{bmatrix}=\bm A .
 \end{aligned}
\]
两种读法分别是“先沿$\bm e_1$伸长、再旋转”和“先旋转、再沿$\bm e_2$伸长”。伸缩倍数相同，伸缩的方向一个在输入端度量，一个在输出端度量。

与奇异值分解对照，两种顺序的关系更清楚。若$\bm A=\bm U\bm\Sigma\bm V^\top$，则第~\ref{sec:ls-polar-singular-best-orthogonal}节将验证$\bm W=\bm U\bm V^\top$、$\bm H=\bm V\bm\Sigma\bm V^\top$，于是
\[
 \bm K=\bm W\bm H\bm W^\top=\bm U\bm\Sigma\bm U^\top .
\]右伸缩因子沿右奇异向量伸缩，左伸缩因子沿左奇异向量伸缩，倍数都是奇异值。

\section{不可逆矩阵的极分解}
\label{sec:ls-polar-singular-best-orthogonal}

上一节的正交因子由$\bm W=\bm A\bm H^{-1}$求出，这要求$\bm H$可逆，即$\bm A$可逆。奇异方阵的伸缩因子照样由第~\ref{sec:ls-matrix-square-root}节求出，正交因子则要换一种办法。构造分三步：先取奇异值分解$\bm A=\bm U\bm\Sigma\bm V^\top$，令$\bm W=\bm U\bm V^\top$、$\bm H=\bm V\bm\Sigma\bm V^\top$，核验它们构成极分解；再找出正交因子中由$\bm A$确定的部分，即它在$(\ker\bm A)^\perp$上的作用；最后说明其余部分怎样自由选择，得到全部正交因子。

\subsection{由奇异值分解构造}
\label{sec:ls-polar-singular-construction}

当$\bm A$奇异时，$\bm S=\bm A^\top\bm A$仍对称半正定，但有零特征值。定理~\ref{thm:ls-psd-square-root}对半正定矩阵成立，所以沿谱取非负平方根仍能构造唯一的半正定$\bm H=\bm S^{1/2}$。上一节的做法却在下一步失效：零特征值不能取倒数，$\bm H^{-1}$不存在，也就无法由$\bm W=\bm A\bm H^{-1}$得到正交因子。

这时可以直接利用第~\ref{chap:singular-value-decomposition}章的奇异值分解。由定理~\ref{thm:ls-svd-existence}，方阵$\bm A$可以写成
\[
 \bm A=\bm U\bm\Sigma\bm V^\top ,
\]
其中$\bm U,\bm V\in\mathbb R^{d\times d}$正交，$\bm\Sigma=\operatorname{diag}(\sigma_1,\ldots,\sigma_d)$，$\sigma_1\geq\cdots\geq\sigma_d\geq0$。在$\bm U$与$\bm\Sigma$之间插入$\bm V^\top\bm V=\bm I$，就得到候选
\begin{equation}
\begin{gathered}
 \bm A=(\bm U\bm V^\top)(\bm V\bm\Sigma\bm V^\top)\\
 \bm W=\bm U\bm V^\top,\quad \bm H=\bm V\bm\Sigma\bm V^\top .
\end{gathered}
\label{eq:ls-polar-from-svd}
\end{equation}

\begin{theorem}[一般方阵的极分解]
\label{thm:ls-polar-general}
每个$\bm A\in\mathbb R^{d\times d}$都有极分解，式~\eqref{eq:ls-polar-from-svd}给出其中一个。伸缩因子$\bm H=(\bm A^\top\bm A)^{1/2}$总是唯一的；正交因子$\bm W$唯一，当且仅当$\bm A$可逆。
\end{theorem}
\begin{proof}
先验证式~\eqref{eq:ls-polar-from-svd}。$\bm H=\bm V\bm\Sigma\bm V^\top$是对称矩阵的正交谱分解，特征值$\sigma_i\geq0$，故$\bm H$对称半正定；由$\bm U^\top\bm U=\bm I$与$\bm\Sigma^\top=\bm\Sigma$，
\[
 \bm H^2=\bm V\bm\Sigma^2\bm V^\top
 =\bm V\bm\Sigma^\top\bm U^\top\bm U\bm\Sigma\bm V^\top
 =\bm A^\top\bm A .
\]
又
\[
 (\bm U\bm V^\top)^\top\bm U\bm V^\top=\bm V\bm U^\top\bm U\bm V^\top=\bm I ,
\]
所以$\bm W$正交；而
\[
 \bm W\bm H=\bm U\bm V^\top\bm V\bm\Sigma\bm V^\top=\bm U\bm\Sigma\bm V^\top=\bm A .
\]
这就证明了存在性。

$\bm H$的唯一性与可逆情形相同：任何极分解都满足式~\eqref{eq:ls-polar-square-root-equation}，由定理~\ref{thm:ls-psd-square-root}，$\bm H=(\bm A^\top\bm A)^{1/2}$。$\bm A$可逆时，$\bm W$的唯一性已由定理~\ref{thm:ls-polar-invertible}给出。

设$\bm A$奇异，则$\operatorname{rank}\bm A<d$，由定理~\ref{thm:ls-svd-existence}知$\sigma_d=0$。令$\bm U'$为把$\bm U$的最后一列$\bm u_d$换成$-\bm u_d$所得的矩阵，它仍正交。$\bm U\bm\Sigma$的第$i$列是$\sigma_i\bm u_i$，最后一列因$\sigma_d=0$而为零，所以$\bm U'\bm\Sigma=\bm U\bm\Sigma$，从而
\[
 \bm A=\bm U'\bm\Sigma\bm V^\top .
\]
按同样的验证，$\bm W'=\bm U'\bm V^\top$与同一个$\bm H$也构成极分解，而
\[
 \bm W-\bm W'=(\bm U-\bm U')\bm V^\top=2\bm u_d\bm v_d^\top\neq\bm0 ,
\]
所以正交因子不唯一。
\end{proof}

可逆时，式~\eqref{eq:ls-polar-from-svd}与定理~\ref{thm:ls-polar-invertible}给出的是同一对因子，因为两者都是极分解，而可逆情形的极分解唯一。谱开方的路线解释了因子为什么是这个样子；实际计算时，若已经有SVD，可以直接读出两个因子。

\subsection{规范部分等距因子}
\label{sec:ls-polar-partial-isometry}

不唯一的原因可以从作用上看清。设$r=\operatorname{rank}\bm A$。对$i\leq r$，$\bm H\bm v_i=\sigma_i\bm v_i$且$\sigma_i>0$，由$\bm A=\bm W\bm H$得
\[
 \bm W\bm v_i=\bm A\bm v_i/\sigma_i=\bm u_i ,
\]
正交因子在这些方向上由$\bm A$完全确定。对$i>r$，$\bm H$把$\bm v_i$压成零，$\bm A$在这些方向上不提供任何信息。

所以分解中真正由$\bm A$确定的，是正交因子在$\operatorname{span}\{\bm v_1,\ldots,\bm v_r\}=(\ker\bm A)^\perp$上的作用。把这部分单独取出，在$\ker\bm A$上规定为零，得到
\begin{equation}
 \bm W_0=\bm U_r\bm V_r^\top=\sum_{i=1}^r\bm u_i\bm v_i^\top ,
 \label{eq:ls-polar-partial-isometry}
\end{equation}
其中$\bm U_r$、$\bm V_r$是秩分解形式中的因子。它把$\bm v_i$送到$\bm u_i$（$i\leq r$），把$\ker\bm A$送到零。对$(\ker\bm A)^\perp$中的$\bm x=\sum_{i\leq r}c_i\bm v_i$，
\[
 \lVert\bm W_0\bm x\rVert_2^2=\sum_{i\leq r}c_i^2=\lVert\bm x\rVert_2^2 .
\]

一般地，若矩阵在其核的正交补上保持长度，就称它为\term{部分等距}（partial isometry）。$\bm W_0$不是正交矩阵，但它在需要保持长度的地方保持长度。

$\bm W_0$由$\bm A$唯一确定，不依赖所选的奇异值分解。由$\bm H=\bm V_r\bm\Sigma_r\bm V_r^\top$，
\[
 \bm W_0\bm H=\bm U_r\bm\Sigma_r\bm V_r^\top=\bm A .
\]

反过来，设矩阵$\bm Z$满足$\bm Z\bm H=\bm A$且在$\ker\bm A$上为零。$\operatorname{range}\bm H=\operatorname{span}\{\bm v_1,\ldots,\bm v_r\}$，其中每个$\bm y=\bm H\bm x$都满足$\bm Z\bm y=\bm A\bm x$，所以$\bm Z$在$\operatorname{range}\bm H$上由$\bm A$决定；它在正交补$\ker\bm A$上又规定为零，所以$\bm Z$处处确定，即$\bm Z=\bm W_0$。称$\bm W_0$为$\bm A$的规范部分等距因子。

\subsection{全部正交因子}
\label{sec:ls-polar-all-orthogonal-factors}

全部正交因子都由$\bm W_0$补全得到。记$\bm V_0=[\bm v_{r+1},\ldots,\bm v_d]$、$\bm U_0=[\bm u_{r+1},\ldots,\bm u_d]$，它们分别是$\ker\bm A$与$\operatorname{range}(\bm A)^\perp$的标准正交基。

\begin{theorem}[奇异方阵的全部正交因子]
\label{thm:polar-singular-orthogonal-factors}
设$\bm A\in\mathbb R^{d\times d}$，$r=\operatorname{rank}\bm A<d$，记号如上。则正交矩阵$\bm W$满足$\bm A=\bm W(\bm A^\top\bm A)^{1/2}$，当且仅当存在$d-r$阶正交矩阵$\bm R$，使
\begin{equation}
 \bm W=\bm W_0+\bm U_0\bm R\bm V_0^\top .
 \label{eq:ls-polar-all-orthogonal-factors}
\end{equation}
\end{theorem}
\begin{proof}
充分性：$\bm U_0\bm R\bm V_0^\top$在$(\ker\bm A)^\perp$上为零，所以$\bm W\bm H=\bm W_0\bm H=\bm A$。又
\[
 \bm W[\bm V_r,\bm V_0]=[\bm U_r,\bm U_0\bm R],
\]
右端各列标准正交：$\bm U_r$与$\bm U_0$的列分别标准正交且彼此正交，$\bm R$正交。所以$\bm W$把标准正交基送到标准正交基，是正交矩阵。

必要性：设$\bm W$正交且$\bm W\bm H=\bm A$。由前面的分析，$\bm W\bm V_r=\bm U_r$。正交矩阵保持内积，所以$\bm W\bm V_0$的列标准正交，并且与$\bm W\bm V_r=\bm U_r$的列正交，即落在$\operatorname{range}(\bm A)^\perp$中。于是
\[
 \bm W\bm V_0=\bm U_0\bm R,\qquad\bm R=\bm U_0^\top\bm W\bm V_0 .
\]

再验证$\bm R$正交：
\[
 \begin{aligned}
 \bm R^\top\bm R&=\bm V_0^\top\bm W^\top\bm U_0\bm U_0^\top\bm W\bm V_0\\
 &=\bm V_0^\top\bm W^\top\bm W\bm V_0=\bm I ,
 \end{aligned}
\]
这里用到$\bm U_0\bm U_0^\top$是到$\operatorname{range}(\bm A)^\perp$的正交投影，它不改变$\bm W\bm V_0$的列。$\bm R$是方阵，所以正交。

最后，
\[
 \begin{aligned}
 \bm W&=\bm W[\bm V_r,\bm V_0][\bm V_r,\bm V_0]^\top\\
 &=\bm U_r\bm V_r^\top+\bm U_0\bm R\bm V_0^\top .
 \end{aligned}
\]
\end{proof}

定理说明，正交因子的全部自由度都来自核：它必须把$\ker\bm A$保长地送到$\operatorname{range}(\bm A)^\perp$，至于怎样送，$\bm A$没有提供任何信息。核的维数$d-r$越大，可选的$\bm R$越多；可逆时$d-r=0$，$\bm W=\bm W_0$唯一。

定理~\ref{thm:ls-polar-general}证明中的变号只是取$\bm R=-1$（$d-r=1$时）或$\bm R=\operatorname{diag}(1,\ldots,1,-1)$的特例。这里始终没有定义不存在的$\bm H^{-1}$，也不能用逐元素取倒数代替矩阵求逆。

\begin{example}[秩一方阵的极分解]
\label{ex:polar-rank-one}
用最简单的奇异方阵看一看定理中的自由度。下面用以下矩阵演示：
\[
 \bm A=\begin{bmatrix}0&0\\1&0\end{bmatrix}.
\]
它把$\bm e_1$送到$\bm e_2$，把$\bm e_2$压成零。伸缩因子照常求出：
\[
 \bm A^\top\bm A=\operatorname{diag}(1,0),\qquad\bm H=\operatorname{diag}(1,0).
\]

正交因子在$\bm e_1$方向上被$\bm A$确定：取$\bm v_1=\bm e_1$、$\bm u_1=\bm e_2$，规范部分等距因子为
\[
 \bm W_0=\bm e_2\bm e_1^\top=\begin{bmatrix}0&0\\1&0\end{bmatrix},
\]
恰好就是$\bm A$本身。

剩下的是怎样把核$\ker\bm A=\operatorname{span}\{\bm e_2\}$保长地送到$\operatorname{range}(\bm A)^\perp=\operatorname{span}\{\bm e_1\}$。取$\bm V_0=\bm e_2$、$\bm U_0=\bm e_1$，这里$d-r=1$，$\bm R$只能是$\pm1$，于是得到两个正交因子
\[
 \bm W_\pm=\begin{bmatrix}0&\pm1\\1&0\end{bmatrix}.
\]
两者都满足$\bm W_\pm\bm H=\bm A$，因为$\bm H$的第二列为零，$\bm W_\pm$第二列的差别被乘掉了。

从几何上看，$\bm W_+$是关于直线$x_1=x_2$的反射，$\bm W_-$是逆时针旋转$90$度，类型都不同；它们只在$\bm e_1$上必须一致，在被压成零的$\bm e_2$方向上，$\bm A$不提供任何约束。
\end{example}

\section{矩形矩阵的极分解}
\label{sec:ls-polar-rectangular}

矩形矩阵$\bm A\in\mathbb R^{m\times d}$的输入与输出不在同一空间，不存在把$\mathbb R^d$映到$\mathbb R^m$的正交矩阵，保长因子只能是列标准正交或行标准正交的矩阵。哪一种可行，取决于哪一端更大。构造的办法是：高矩阵取薄奇异值分解，仿照方阵写出列标准正交的保长因子与伸缩因子，满列秩时也可以像可逆方阵那样直接撤销伸缩；宽矩阵则对转置使用高矩阵的结论，再转置回来，得到先保长、后伸缩的左极分解。

\subsection{高矩阵}
\label{sec:ls-polar-tall}

设$m\geq d$，即$\bm A$是高矩阵。取薄奇异值分解$\bm A=\bm U_d\bm\Sigma_d\bm V^\top$，其中$\bm U_d\in\mathbb R^{m\times d}$列标准正交，$\bm\Sigma_d=\operatorname{diag}(\sigma_1,\ldots,\sigma_d)$，$\bm V\in\mathbb R^{d\times d}$正交。仿照式~\eqref{eq:ls-polar-from-svd}，令
\[
 \bm W=\bm U_d\bm V^\top\in\mathbb R^{m\times d},\qquad
 \bm H=\bm V\bm\Sigma_d\bm V^\top\in\mathbb R^{d\times d}.
\]
则
\[
 \bm W^\top\bm W=\bm V\bm U_d^\top\bm U_d\bm V^\top=\bm I_d ,
\]
即$\bm W$的列标准正交，它把$\mathbb R^d$保长地映入$\mathbb R^m$。又
\[
 \bm W\bm H=\bm U_d\bm\Sigma_d\bm V^\top=\bm A .
\]

唯一性的讨论与方阵相同。由$\bm W^\top\bm W=\bm I_d$，仍有
\[
 \bm A^\top\bm A=\bm H\bm W^\top\bm W\bm H=\bm H^2 ,
\]
所以$\bm H=(\bm A^\top\bm A)^{1/2}$唯一。若$\bm A$满列秩，$\bm H$可逆，$\bm W=\bm A\bm H^{-1}$也唯一；若$\operatorname{rank}\bm A<d$，同样可以在$\ker\bm A$上改变$\bm W$，例如把$\bm U_d$中对应零奇异值的一列变号。

满列秩时不必先求奇异值分解。这时$\bm A^\top\bm A$正定，与可逆方阵完全一样，$\bm H=(\bm A^\top\bm A)^{1/2}$可逆，
\[
 \bm W=\bm A(\bm A^\top\bm A)^{-1/2}.
\]
直接验证：
\[
 \bm W^\top\bm W=\bm H^{-1}\bm A^\top\bm A\bm H^{-1}=\bm I_d,\qquad\bm W\bm H=\bm A .
\]

\subsection{宽矩阵}
\label{sec:ls-polar-wide}

设$m<d$，即$\bm A$是宽矩阵。这时不存在列标准正交的$m\times d$矩阵，因为$\mathbb R^m$中至多有$m$个标准正交向量。改对$\bm A^\top$用上面的结论：
\[
 \bm A^\top=\bm W_1\bm H_1 ,
\]
其中$\bm W_1\in\mathbb R^{d\times m}$列标准正交，$\bm H_1=(\bm A\bm A^\top)^{1/2}$。取转置，
\[
 \bm A=\bm H_1\bm W_1^\top ,
\]
其中$\bm W_1^\top\in\mathbb R^{m\times d}$行标准正交。这是左极分解的形式：先把$\mathbb R^d$正交地投到一个$m$维坐标系中，再在$\mathbb R^m$中伸缩。所以对宽矩阵，自然的是左极分解；对高矩阵，自然的是右极分解。

\begin{example}[行向量与列向量的极分解]
\label{ex:polar-rectangular}
最小的矩形矩阵是单个列向量或行向量，两种情形正好对应高矩阵和宽矩阵。先看列向量$\bm a=(3,4)^\top\in\mathbb R^{2\times1}$，它把一个数映成平面中的向量。$\bm a^\top\bm a=25$，所以
\[
 \bm H=5,\qquad\bm W=\bm a/5=(0.6,0.8)^\top,\qquad\bm a=\bm W\bm H .
\]
这个右极分解读作：先把输入的数伸长$5$倍，再保长地放到方向$(0.6,0.8)^\top$上。

转置$\bm a^\top=[3,4]$是宽矩阵，把平面向量映成一个数，这时应取左极分解。左伸缩因子来自$\bm a^\top(\bm a^\top)^\top=\bm a^\top\bm a=25$，开方得$5$，于是
\[
 \bm a^\top=5\cdot[0.6,0.8].
\]
行标准正交的$[0.6,0.8]$先读出输入在单位方向$(0.6,0.8)^\top$上的坐标，再乘以$5$；保长的一步在前，伸缩在后，顺序与列向量的情形相反。
\end{example}

\section{矩阵函数的谱表示}
\label{sec:ls-matrix-functions}
\label{sec:matrix-analysis-functions-structured-operations}

前面几节反复把一个标量运算沿特征方向作用到矩阵上。本节把这种做法写成一般定义，检查它与矩阵运算相容，并指出它的适用范围。

\paragraph{沿谱运算的共同方法}

回看前面两次构造。求$\bm H$时，我们保持$\bm S$的特征方向，把每个$\lambda_i$换成$\sqrt{\lambda_i}$；求$\bm H^{-1}$时，同样保持特征方向，把$\lambda_i$换成$1/\sqrt{\lambda_i}$。两次都是“方向不动，只对特征值作用一个标量函数”。把这个实际用过的操作归纳出来，就得到对称矩阵的函数。

设实对称矩阵$\bm S$有正交谱分解$\bm S=\bm Q\bm\Lambda\bm Q^\top$。按第~\ref{sec:ls-spectral-projections}节，把同一特征值的项合并，可以写成
\[
 \bm S=\sum_\lambda\lambda\bm P_\lambda ,
\]
求和遍历$\bm S$的不同特征值，$\bm P_\lambda$是到特征空间$\ker(\bm S-\lambda\bm I)$的正交投影，满足
\[
 \bm P_\lambda\bm P_\mu=\bm0\ (\lambda\neq\mu),\qquad\bm P_\lambda^2=\bm P_\lambda,\qquad\sum_\lambda\bm P_\lambda=\bm I .
\]

\begin{definition}[对称矩阵的谱函数]
\label{def:ls-symmetric-matrix-function}
设$\bm S\in\mathbb R^{d\times d}$实对称，$f$是定义在$\bm S$全部特征值上的实函数。令
\begin{equation}
\begin{aligned}
 f(\bm S)
 &=\bm Q\operatorname{diag}\bigl(f(\lambda_1),\ldots,f(\lambda_d)\bigr)\bm Q^\top\\
 &=\sum_\lambda f(\lambda)\bm P_\lambda ,
\end{aligned}
\label{eq:ls-symmetric-spectral-function}
\end{equation}
称为由$f$沿谱定义的\term{矩阵函数}（matrix function）。
\end{definition}

两种写法相等：按列展开，第一式是$\sum_if(\lambda_i)\bm q_i\bm q_i^\top$；把$\lambda_i=\lambda$的项合并，
\[
 \sum_{\lambda_i=\lambda}\bm q_i\bm q_i^\top=\bm P_\lambda ,
\]
就得到第二式。

第二式只涉及特征值和谱投影，而谱投影由$\bm S$唯一确定，所以$f(\bm S)$与重根空间内基的选取无关，定义是良好的。

$f(z)=\sqrt z$时，它回到第~\ref{sec:ls-matrix-square-root}节的$\bm S^{1/2}$；$f(z)=1/\sqrt z$时，回到第~\ref{sec:ls-polar-orthogonal-factor}节的$\bm S^{-1/2}$。前面那些构造在这里得到总结，而不是反过来，必须先建立一般的矩阵函数才能解出$\bm H$与$\bm W$。

推广之后要检查它与已有的矩阵运算相容，否则同一个记号会有两种含义。

\begin{lemma}[谱函数与矩阵运算相容]
\label{lem:ls-spectral-function-calculus}
设$\bm S$实对称，$f,g$定义在$\bm S$的全部特征值上。
\begin{enumerate}
 \item 对实系数多项式$p(z)=\sum_{k=0}^mc_kz^k$，按矩阵幂与加法算出的$\sum_kc_k\bm S^k$等于式~\eqref{eq:ls-symmetric-spectral-function}定义的$p(\bm S)$。
 \item $f(\bm S)+g(\bm S)=(f+g)(\bm S)$，$f(\bm S)g(\bm S)=(fg)(\bm S)$。
 \item 若$f$在全部特征值上不为零，则$f(\bm S)$可逆，且$f(\bm S)^{-1}=(1/f)(\bm S)$；特别地，$\bm S$可逆时取$f(z)=1/z$即得$\bm S^{-1}$。
\end{enumerate}
\end{lemma}
\begin{proof}
先证第二条。加法由式~\eqref{eq:ls-symmetric-spectral-function}的第二式直接得到。乘法利用谱投影的正交关系：
\[
 f(\bm S)g(\bm S)
 =\sum_\lambda\sum_\mu f(\lambda)g(\mu)\bm P_\lambda\bm P_\mu
 =\sum_\lambda f(\lambda)g(\lambda)\bm P_\lambda
 =(fg)(\bm S),
\]
中间一步用到$\lambda\neq\mu$时$\bm P_\lambda\bm P_\mu=\bm0$，以及$\bm P_\lambda^2=\bm P_\lambda$。

第一条：常数函数$1$给出$\sum_\lambda\bm P_\lambda=\bm I=\bm S^0$，恒等函数$z$给出$\sum_\lambda\lambda\bm P_\lambda=\bm S$。由乘法规则对$k$归纳，$z^k$给出$\bm S^k$；再由加法规则（数乘同理）按系数相加，得到$p(\bm S)=\sum_kc_k\bm S^k$。

第三条：由乘法规则，$f(\bm S)\,(1/f)(\bm S)=(f\cdot 1/f)(\bm S)$，即常数函数$1$作用的结果$\bm I$；两个因子的顺序交换后同理。$f(z)=1/z$时，$(1/f)(z)=z$，于是$(1/z)(\bm S)$是$\bm S$的逆。
\end{proof}

第一条也可以直接看出：连乘$\bm S^k=\bm Q\bm\Lambda\bm Q^\top\cdots\bm Q\bm\Lambda\bm Q^\top$时相邻的$\bm Q^\top\bm Q$相消，得
\[
 \bm S^k=\bm Q\bm\Lambda^k\bm Q^\top .
\]
所以对多项式而言，沿谱作用与按矩阵运算计算是一回事；定义~\ref{def:ls-symmetric-matrix-function}只是把这个事实推广到一般的$f$。

由第二条，$(\sqrt z)^2=z$给出$(\bm S^{1/2})^2=\bm S$，$\sqrt z\cdot(1/\sqrt z)=1$给出$\bm S^{1/2}\bm S^{-1/2}=\bm I$，与前面的直接验证一致。

矩阵函数沿特征方向作用，而不是对每个元素作用。下面的例子对照这两种做法。

\begin{example}[矩阵平方根与逐元素开方]
\label{ex:ls-matrix-square-root-two-by-two}
仍取示例~\ref{ex:ls-symmetric-spectral-two-by-two}中的矩阵及其谱分解：
\[
 \bm S=\begin{bmatrix}2&1\\1&2\end{bmatrix}=3\bm q_1\bm q_1^\top+\bm q_2\bm q_2^\top ,
\]
其中$\bm q_1=(1,1)^\top/\sqrt2$，$\bm q_2=(1,-1)^\top/\sqrt2$。

按式~\eqref{eq:ls-symmetric-spectral-function}取$f(z)=\sqrt z$，保持两个投影不变，只把系数$3$与$1$换成平方根：
\[
 \bm S^{1/2}
 =\sqrt3\,\bm q_1\bm q_1^\top+\bm q_2\bm q_2^\top
 =\frac12\begin{bmatrix}\sqrt3+1&\sqrt3-1\\\sqrt3-1&\sqrt3+1\end{bmatrix}.
\]
为了与后面的逐元素做法比较，直接计算它的平方。对角元与非对角元分别为
\[
 \begin{aligned}
 &\tfrac14\bigl[(\sqrt3+1)^2+(\sqrt3-1)^2\bigr]=\tfrac14\cdot8=2,\\
 &\tfrac14\cdot2(\sqrt3+1)(\sqrt3-1)=\tfrac14\cdot4=1,
 \end{aligned}
\]
所以$(\bm S^{1/2})^2=\bm S$。

同样取$f(z)=1/\sqrt z$，得逆平方根
\[
 \bm S^{-1/2}
 =\frac1{\sqrt3}\bm q_1\bm q_1^\top+\bm q_2\bm q_2^\top
 =\frac1{2\sqrt3}\begin{bmatrix}\sqrt3+1&1-\sqrt3\\1-\sqrt3&\sqrt3+1\end{bmatrix}.
\]
若改为对每个元素开方，得到的矩阵平方以后是
\[
 \begin{bmatrix}\sqrt2&1\\1&\sqrt2\end{bmatrix}^2
 =\begin{bmatrix}3&2\sqrt2\\2\sqrt2&3\end{bmatrix}\neq\bm S .
\]
差别在于矩阵乘法会把不同元素混合起来：平方的每个元素是一行与一列的内积，而不是单个元素的平方。沿特征方向作用时，各方向在乘法下互不混合，开方才与平方相容；逐元素开方没有这种结构，所以得不到$\bm S$。
\end{example}

固定的矩阵只有有限个特征值，所以式~\eqref{eq:ls-symmetric-spectral-function}只需要$f$在这些点上的值，不要求$f$连续，也不需要矩阵级数或微分。

还要把沿谱构造的矩阵函数与第~\ref{chap:matrix-differentiation}章研究的一般矩阵输入映射区分开。行列式$\bm X\mapsto\det\bm X$、矩阵乘法$(\bm X,\bm Y)\mapsto\bm X\bm Y$、求逆等都以矩阵为输入，范围比$\bm S\mapsto f(\bm S)$广得多。

即使研究$\bm S\mapsto f(\bm S)$本身的微分，扰动$\bm E$一般也不保持$\bm S$的特征方向，不能直接套用标量求导公式。例如$f(z)=z^2$，
\[
 (\bm S+\bm E)^2-\bm S^2=\bm S\bm E+\bm E\bm S+\bm E^2 ,
\]
一阶部分是$\bm S\bm E+\bm E\bm S$；只有$\bm S$与$\bm E$可交换时，它才化为标量公式对应的$2\bm S\bm E$。

\paragraph{在已知分解上作用谱函数}

谱函数的定义只要求一组正交的特征方向，所以手中已有的分解可以直接使用。若$\bm S=\bm Q\bm\Lambda\bm Q^\top$，而$\bm Z$是正交矩阵，则$\bm Z\bm S\bm Z^\top=(\bm Z\bm Q)\bm\Lambda(\bm Z\bm Q)^\top$也是正交谱分解，于是
\[
 f(\bm Z\bm S\bm Z^\top)=\bm Z f(\bm S)\bm Z^\top .
\]
也就是说，先用正交矩阵转动、再作用谱函数，与先作用谱函数、再转动，结果相同。第~\ref{sec:ls-polar-left}节的左伸缩因子也可以这样得到。由
\[
 \bm A\bm A^\top=\bm W\bm H\bm H\bm W^\top=\bm W(\bm A^\top\bm A)\bm W^\top ,
\]
取$f(z)=\sqrt z$得
\[
 \bm K=(\bm A\bm A^\top)^{1/2}=\bm W(\bm A^\top\bm A)^{1/2}\bm W^\top=\bm W\bm H\bm W^\top .
\]

奇异值分解同样提供现成的谱分解。若$\bm A=\bm U\bm\Sigma\bm V^\top$，则
\[
 \bm A^\top\bm A=\bm V(\bm\Sigma^\top\bm\Sigma)\bm V^\top,\qquad\bm A\bm A^\top=\bm U(\bm\Sigma\bm\Sigma^\top)\bm U^\top ,
\]
所以
\[
 \begin{gathered}
 f(\bm A^\top\bm A)=\bm V f(\bm\Sigma^\top\bm\Sigma)\bm V^\top ,\\
 f(\bm A\bm A^\top)=\bm U f(\bm\Sigma\bm\Sigma^\top)\bm U^\top ,
 \end{gathered}
\]
其中$f$作用于对角矩阵的各个对角元。

取$f(z)=\sqrt z$，就得到$\bm H=\bm V\bm\Sigma\bm V^\top$与$\bm K=\bm U\bm\Sigma\bm U^\top$。可逆方阵再取$f(z)=1/\sqrt z$，得
\[
 \begin{aligned}
 \bm W=\bm A(\bm A^\top\bm A)^{-1/2}
 &=\bm U\bm\Sigma\bm V^\top\bm V\bm\Sigma^{-1}\bm V^\top\\
 &=\bm U\bm V^\top ,
 \end{aligned}
\]
与式~\eqref{eq:ls-polar-from-svd}一致。

\paragraph{实幂与对数的补充例子}

对正定$\bm S$，全部$\lambda_i>0$，任意实数$\alpha$都使$z^\alpha$在谱上有定义，于是
\[
 \bm S^\alpha=\bm Q\operatorname{diag}(\lambda_1^\alpha,\ldots,\lambda_d^\alpha)\bm Q^\top .
\]
由引理~\ref{lem:ls-spectral-function-calculus}，
\[
 \bm S^\alpha\bm S^\beta=\bm S^{\alpha+\beta}.
\]

特别地，对正整数$m$，$(\bm S^{1/m})^m=\bm S$，即把非负伸缩分成$m$次相同的作用时，每次沿各主方向伸缩$\lambda_i^{1/m}$倍；$\bm S^{1/m}$本身对称正定，因为它的谱分解中特征值$\lambda_i^{1/m}>0$。这里只用到它是$\bm X^m=\bm S$的一个对称正定解，本节不讨论$m>2$时这类解的唯一性。

零特征值限制了可用的幂。若$\bm S$半正定但有零特征值，$\alpha>0$时$0^\alpha=0$，$\bm S^\alpha$仍有实值的谱定义；$\alpha<0$时$0^\alpha$没有定义，负幂不能沿用。前面平方根只要求半正定、逆平方根却要求正定，正是这一区别。负特征值则使平方根、对数等运算失去实值定义，例如$\sqrt{-1}$不是实数，所以本节的实幂与对数只对半正定或正定矩阵使用。

对正定$\bm S$还可以取$f(z)=\log z$，得到\term{矩阵对数}（matrix logarithm）
\[
 \log\bm S=\bm Q\operatorname{diag}(\log\lambda_1,\ldots,\log\lambda_d)\bm Q^\top .
\]
它记录的是各主方向上伸缩倍数的对数：$\log\lambda_i>0$表示沿$\bm q_i$拉长，$\log\lambda_i<0$表示压缩。上式本身就是对称矩阵$\log\bm S$的正交谱分解，特征值为$\log\lambda_i$。

对数与指数互为逆映射。由定义~\ref{def:ls-symmetric-matrix-function}与基的选取无关，可以用这组$\bm Q$对$\log\bm S$作用$f(z)=e^z$（矩阵指数在本节也按式~\eqref{eq:ls-symmetric-spectral-function}沿谱定义），得
\[
 e^{\log\bm S}=\bm Q\operatorname{diag}(e^{\log\lambda_i})\bm Q^\top=\bm S .
\]
反过来，对任意实对称$\bm T=\bm Q\operatorname{diag}(\mu_i)\bm Q^\top$，
\[
 e^{\bm T}=\bm Q\operatorname{diag}(e^{\mu_i})\bm Q^\top
\]
是特征值全为正的对称矩阵，即正定，同理$\log e^{\bm T}=\bm T$。

所以在沿谱的定义下，$\log$与$e^{(\cdot)}$在正定矩阵与实对称矩阵之间互为逆映射。$\bm S$有零特征值时$\log0$无定义，这一定义不能使用。

要区分矩阵$\log\bm S$与标量$\log\det\bm S$。由$\det\bm S=\det\bm\Lambda=\prod_i\lambda_i$及迹的循环性，
\[
 \begin{aligned}
 \log\det\bm S
 &=\sum_{i=1}^d\log\lambda_i\\
 &=\operatorname{tr}(\bm Q\operatorname{diag}(\log\lambda_i)\bm Q^\top)
 =\operatorname{tr}(\log\bm S).
 \end{aligned}
\]
左端是一个数，$\log\bm S$是$d$阶矩阵，二者通过迹联系。对数在本节只是谱函数的又一个例子，不参与极分解的证明。第~\ref{sec:multivariable-logdet-differential}节求$\log\det$的微分时，使用的是行列式的微分与标量对数的链式法则，并不以矩阵对数的微分为前提。

\FloatBarrier
\section{本章小结}
\label{sec:polar-chapter-summary}

本章开头问：能否把一个线性作用分成沿正交方向的伸缩与保长的转向？对可逆方阵，答案完全确定。由正交因子保持长度，$\bm A=\bm W\bm H$迫使$\bm H^2=\bm A^\top\bm A$；沿$\bm A^\top\bm A$的谱取正平方根，得到唯一的正定伸缩因子$\bm H=(\bm A^\top\bm A)^{1/2}$，再由$\bm W=\bm A\bm H^{-1}$得到唯一的正交因子。$\bm W$的行列式与$\bm A$同号，所以它可能包含反射。

伸缩与转向的顺序可以交换。左极分解$\bm A=\bm K\bm W$使用同一个正交因子，伸缩因子换成$\bm K=(\bm A\bm A^\top)^{1/2}=\bm W\bm H\bm W^\top$：伸缩倍数不变，伸缩方向从输入端转到输出端。用奇异值分解表示，$\bm W=\bm U\bm V^\top$、$\bm H=\bm V\bm\Sigma\bm V^\top$、$\bm K=\bm U\bm\Sigma\bm U^\top$。

奇异方阵仍有极分解，伸缩因子仍唯一，正交因子则不唯一。由$\bm A$确定的只有规范部分等距因子$\bm W_0=\bm U_r\bm V_r^\top$，即正交因子在$(\ker\bm A)^\perp$上的作用；全部正交因子是$\bm W_0+\bm U_0\bm R\bm V_0^\top$，自由度来自怎样把核保长地送到值域的正交补。矩形矩阵的保长因子只能是列标准正交（高矩阵，右极分解）或行标准正交（宽矩阵，左极分解）的矩阵。

开方与求逆都是“方向不动，只改变特征值”，由此得到对称矩阵的谱函数$f(\bm S)=\sum_\lambda f(\lambda)\bm P_\lambda$。它与基的选取无关，与多项式运算、乘法和求逆相容，只需$f$在特征值上有定义；零特征值排除了负幂、逆平方根与对数，负特征值排除了实平方根与对数。已有的谱分解或奇异值分解可以直接用来计算这些函数。

极分解的正交因子在Frobenius范数下是离$\bm A$最近的正交矩阵，这一最优性在第~\ref{chap:matrix-approximation}章证明。随后转向误差：第~\ref{chap:condition-and-spectral-perturbation}章用奇异值度量输入误差的放大，并估计特征值、奇异值与特征子空间受矩阵扰动时的变化。
